% stephane [DOT] adjemian [AT] univ [DASH] lemans [DOT] fr
\documentclass[10pt,a4paper,notitlepage,twocolumn]{article}
\usepackage{amsmath}
\usepackage{amssymb}
\usepackage{amsbsy}
\usepackage{bm}
\usepackage[T1]{fontenc}
\usepackage[utf8x]{inputenc}
\usepackage{palatino}
\usepackage{scrtime}
\usepackage[french]{babel}
\usepackage{float}
\usepackage{nicefrac}

%%%%%%%%%%%%%%%%%%%%%%%%%%%%%%%%%%%%%%%%%%%%%%%%%%%%%%%%%%%%%%%%%%%%%%%%%%%%%%%%%%%%%%%%%%%%%%%%%%%%%

\newcommand{\exercice}[1]{\textsc{\textbf{Exercice}} #1}
\newcommand{\question}[1]{\textbf{(#1)}}
\setlength{\parindent}{0cm}



\begin{document}


\title{\textsc{Séries temporelles}}
\author{\textsc{Université du Mans (2\textsuperscript{nde} session, L3)}}
\date{Mardi 16 juin 2026}

\pagenumbering{gobble}

\maketitle

\exercice{1} Soit le processus moyenne mobile MA(2)~:
\[
  Y_t = \varepsilon_t - \frac{5}{6}\varepsilon_{t-1} + \frac{1}{6}\varepsilon_{t-2}
\]
où $\varepsilon_t$ est un bruit blanc d'espérance nulle et de variance
$\sigma_{\varepsilon}^2$.\newline

\question{1} Ce processus est-il stationnaire au second ordre~?
Justifier sans calcul.\newline

\question{2} Calculer l'espérance et les autocovariances $\gamma(0)$,
$\gamma(1)$, $\gamma(2)$ ainsi que $\gamma(h)$ pour $h\geq 3$.\newline

\question{3} En déduire la fonction d'autocorrélation $\rho(h)$.\newline

\question{4} Le processus est-il inversible~? On déterminera les
racines du polynôme retard moyenne mobile et on donnera sa
factorisation.\newline

\bigskip

\exercice{2} Soit le processus AR(2)~:
\[
  Y_t = 1 + Y_{t-1} - \frac{1}{2}Y_{t-2} + \varepsilon_t
\]
avec $\varepsilon_t$ un bruit blanc d'espérance nulle et de variance
$\sigma_{\varepsilon}^2$.\newline

\question{1} Montrer que ce processus est asymptotiquement
stationnaire au second ordre. Que remarque-t-on sur les racines du
polynôme retard~?\newline

On suppose dans toute la suite que le processus est stationnaire au
second ordre.\newline

\question{2} Calculer son espérance.\newline

\question{3} Écrire les équations de Yule-Walker et en déduire
$\gamma(0)$, $\gamma(1)$ et $\gamma(2)$ en fonction de
$\sigma_{\varepsilon}^2$.\newline

\question{4} Donner la relation de récurrence vérifiée par $\gamma(h)$
pour $h\geq 1$. En résolvant cette récurrence, donner la forme générale
de $\gamma(h)$ et commenter le comportement de la fonction
d'autocorrélation lorsque $h$ varie.\newline

\bigskip

\exercice{3} Soit le processus AR(1) gaussien~:
\[
  y_t = c + \varphi y_{t-1} + \varepsilon_t,\qquad |\varphi|<1
\]
avec $\varepsilon_t \sim \mathcal{N}(0,\sigma_{\varepsilon}^2)$,
indépendants entre eux. On note $\mathcal Y_T \equiv
\{y_1,y_2,\dots,y_T\}$ l'échantillon observé et
$\bm{\theta}=(c,\varphi,\sigma_{\varepsilon}^2)'$ le vecteur des
paramètres. On supposera le processus stationnaire.\newline

\question{1} À l'aide du théorème de Bayes, écrire la décomposition de
la vraisemblance exacte. On précisera la loi marginale de la première
observation $y_1$.\newline

\question{2} Définir la vraisemblance conditionnelle (on conditionnera
par rapport à $y_1$) et donner l'expression de la log-vraisemblance
conditionnelle.\newline

\question{3} Montrer qu'à $\sigma_{\varepsilon}^2$ fixé, la
maximisation de la log-vraisemblance conditionnelle par rapport à
$(c,\varphi)$ équivaut à un problème de moindres carrés ordinaires. En
déduire l'estimateur sous forme matricielle, puis l'estimateur de
$\sigma_{\varepsilon}^2$.\newline

\question{4} On souhaite tester l'absence de dynamique, c'est-à-dire
$H_0:\varphi=0$ contre $H_1:\varphi\neq 0$. Sachant que
$\sqrt{T}(\hat\varphi-\varphi)\xrightarrow{d}\mathcal
N(0,1-\varphi^2)$, proposer une statistique de test et donner sa loi
asymptotique sous $H_0$. Quelle est la règle de décision au seuil de
5~\%~?\newline

\bigskip

\exercice{4} Soit le processus MA(1)~:
\[
  Y_t = \varepsilon_t + 2\,\varepsilon_{t-1}
\]
où $\varepsilon_t$ est un bruit blanc d'espérance nulle et de variance
$\sigma_{\varepsilon}^2$.\newline

\question{1} Calculer la fonction d'autocovariance puis
l'autocorrélation d'ordre 1 de $\{Y_t\}$.\newline

\question{2} On considère le processus $\tilde Y_t = u_t +
\frac{1}{2}u_{t-1}$ où $u_t$ est un bruit blanc d'espérance nulle et
de variance $\sigma_u^2 = 4\sigma_{\varepsilon}^2$. Montrer que
$\{\tilde Y_t\}$ et $\{Y_t\}$ ont la même fonction d'autocovariance.
Que peut-on en conclure sur l'identification d'un processus MA à partir
de ses moments d'ordre 1 et 2~?\newline

\question{3} Rappeler la définition de l'inversibilité d'un processus
MA. Parmi les deux représentations précédentes, laquelle est
inversible~? Pour cette représentation, exprimer (à l'aide du polynôme
retard) l'innovation en fonction du présent et du passé du processus
observé.\newline

\question{4} Pourquoi impose-t-on la condition d'inversibilité lors de
l'estimation d'un modèle MA~?

\end{document}

%%% Local Variables:
%%% mode: latex
%%% TeX-master: t
%%% End:
